\documentclass[11pt]{article}

\usepackage[a4paper,margin=2.4cm]{geometry}
\usepackage{microtype}
\usepackage{booktabs}
\usepackage{xcolor}
\usepackage{listings}
\usepackage{enumitem}
\usepackage{array}
\usepackage{longtable}
\usepackage[colorlinks=true,linkcolor=blue!50!black,urlcolor=blue!50!black]{hyperref}

\definecolor{kw}{RGB}{20,60,140}
\definecolor{codebg}{RGB}{246,246,242}
\definecolor{errbg}{RGB}{252,244,242}

\lstdefinestyle{gl}{
  basicstyle=\ttfamily\small,
  keywordstyle=\bfseries\color{kw},
  backgroundcolor=\color{codebg},
  frame=single, rulecolor=\color{gray!50},
  framesep=5pt, breaklines=true,
  showstringspaces=false, sensitive=true,
  morekeywords={find,follow,group,count,sum,avg,min,max,compute,return,limit,budget,
    tokens,where,order,by,as,and,or,none,contains,assert,edge,merge,distinct,refine,
    compact,retire,flag,derive,class,keep,upto,any,from,with,into,col,prefer,newest,
    source,reason,continue,after,schema,same,minus,plus,times,over,desc,true,false,via}
}
\lstdefinestyle{err}{
  basicstyle=\ttfamily\small,
  backgroundcolor=\color{errbg},
  frame=single, rulecolor=\color{red!30!gray},
  framesep=5pt, breaklines=true, showstringspaces=false
}

\setlist[itemize]{itemsep=2pt,topsep=4pt}
\setlist[enumerate]{itemsep=2pt,topsep=4pt}

\title{\textbf{theorem v0.2}\\[4pt]
\large Benchmarks, engine work, and what the numbers do and do not support}
\author{}
\date{31 August 2026}

\begin{document}
\maketitle
\vspace{-2.2em}

\begin{center}
\begin{minipage}{0.92\textwidth}
\small
\textbf{Summary.} theorem is a graph query and construction language designed so
that a language model cannot express the mistakes it usually makes. This report
covers a working session that rebuilt its evidence base: four benchmarks (two
regenerated, two new), eight defects found in shipped code, four found in the
benchmark harness, and the engine work that makes the system usable outside a
demo. Every figure is reproducible from per-question data committed to the
repository.

The session began by discovering that the published headline number could not be
reproduced by the code that shipped it. That finding, and the guard that
surfaced it, shape everything that follows.
\end{minipage}
\end{center}

\vspace{1em}
\hrule
\vspace{1em}

\tableofcontents
\newpage

\section{What theorem is, in one page}

Language models write graph queries badly, and they fail in structural,
predictable ways: reversed relationship directions, hallucinated labels,
implicit grouping, long-range bracket errors. Two model generations of scaling
have not fixed it; frontier models still sit near 51--61\% execution accuracy on
public benchmarks.

theorem removes each failure mode by construction rather than by prompting
around it:

\begin{itemize}
  \item \textbf{No direction glyphs.} Edges are traversed by naming the
    \emph{role} you arrive at, not by drawing an arrow. A wrong role is a type
    error caught before execution.
  \item \textbf{One line, one step, one name.} No chaining, no nesting, nothing
    to balance.
  \item \textbf{Explicit staged aggregation.} \texttt{group by} then
    \texttt{count}; adding a column to \texttt{return} can never silently change
    the grouping.
  \item \textbf{Schema-closed vocabulary.} The whole program is verified against
    the live schema before anything runs. Errors name the line, suggest the fix,
    and state \texttt{nothing was executed.}
  \item \textbf{Token-budgeted results} with explicit truncation and
    continuation handles.
\end{itemize}

\begin{lstlisting}[style=gl]
find product where launch_year > 2024 as recent
follow recent uses component as parts
follow parts supplied_by source as sups
group by sups as g
count distinct g.parts as n_parts
return sups.name, n_parts order by n_parts desc
\end{lstlisting}

\noindent The engine is a single-process store: an append-only write-ahead log
plus immutable snapshot runs, with the graph held in memory.

\section{The finding that made the session necessary}

The repository published \textbf{78.02\%} execution accuracy on the full
CypherBench test set. That number was not fabricated, but it did not belong to
the code that shipped.

Frozen query files are named by a hash of the prompt that generated them. The
shipped prompt hashed to \texttt{eb0f4010}; the published run's queries came from
\texttt{3ad56b1c}. The tutorial inside the prompt had been halved after that run
and validated only on the \emph{agent-loop} benchmark, where a repair turn can
hide a regression that one-shot translation cannot.

So the repository could not reproduce its own headline. The guard worked; nobody
had run it. Everything in Section~\ref{sec:results} is a regeneration against the
prompt that actually ships.

\section{Results}
\label{sec:results}

Four benchmarks. Two existed and were regenerated; two did not exist.

\subsection{CypherBench: the full public test set}

All 2{,}348 questions, all seven test graphs, every match category, full
unsampled graphs, zero-shot, one generation, no retry, scored by the benchmark's
own comparator against its published gold answers. The text2cypher row is a
\emph{control} run under the official zero-shot prompt on the official Neo4j
image, not a citation.

\begin{center}
\begin{tabular}{lrrr}
\toprule
 & theorem & text2cypher & delta \\
\midrule
Execution accuracy      & \textbf{77.98\%} & 70.36\% & $+7.6$ \\
Held out (excl.\ \texttt{nba}) & \textbf{76.85\%} & --- & \\
Executable              & 96.6\%  & 95.3\%  & \\
Median execution latency & 0.2\,ms & 67.2\,ms & \\
\bottomrule
\end{tabular}
\end{center}

\noindent Ahead on all seven graphs:

\begin{center}
\begin{tabular}{lrr}
\toprule
graph & theorem & text2cypher \\
\midrule
flight\_accident     & 88.4\% & 76.2\% \\
nba                  & 86.7\% & 71.9\% \\
fictional\_character & 80.0\% & 77.1\% \\
company              & 78.1\% & 70.6\% \\
geography            & 74.6\% & 62.6\% \\
politics             & 73.9\% & 69.0\% \\
movie                & 72.3\% & 68.3\% \\
\bottomrule
\end{tabular}
\end{center}

\texttt{nba} is the graph theorem's prompt was written against, which is why the
held-out figure excluding it is reported alongside. Published 2024 baselines for
context: Claude 3.5 Sonnet 61.6\%, GPT-4o 60.2\%.

Regenerating against the shipped prompt moved the full-set number by
\textbf{0.04 points} (78.02 to 77.98) and left the held-out number unchanged.
Half the prompt for four hundredths of a point. Section~\ref{sec:sampling}
explains why that outcome was not free.

\subsection{The agent loop}

CypherBench measures one-shot translation. An agent writes a query, reads the
error, and tries again. This measures whether it converges, in how many turns,
and for how many tokens, on graphs nothing was tuned on: CypherBench's
\emph{train} split, which shares no schema, question or qid with the test set.
Doubled this session from $n=120$ to $n=240$ across two held-out graphs.

\begin{center}
\begin{tabular}{lrrrrrr}
\toprule
 & $n$ & solve@1 & solve@2 & solve@3 & turns & tokens/q \\
\midrule
theorem     & 240 & \textbf{82.9\%} & 86.7\% & 87.9\% & 1.07 & 2{,}064 \\
text2cypher & 240 & 72.1\% & 82.1\% & 85.0\% & 1.19 & 1{,}068 \\
\bottomrule
\end{tabular}
\end{center}

Of 240 questions, 183 were solved by both and 8 by neither. The verdict rests on
the 49 they disagree about: theorem alone 28, text2cypher alone 21. Exact
McNemar two-sided $p = 0.392$.

\textbf{solve@3 is a tie, and the published page says so.} Reading a winner into
a 2.9-point difference at this sample size would be reading noise.

What doubling the sample surfaced is the \emph{first attempt}: 82.9\% against
72.1\%, a 10.8-point gap that the retries then close, with convergence in 1.07
turns against 1.19. An agent that can retry pays the difference in turns; one
that cannot pays it in answers.

\subsection{Silent failure (new)}

Execution accuracy counts what a model gets right and says nothing about what
happens when it gets one wrong. That is where the two languages differ most, and
nothing measured it.

\textbf{Method.} Take a query known to be correct, break exactly one token the
way models break them, run it, and record what the caller sees. Four mutations
(wrong class, wrong edge type, wrong property, reversed direction), applied to
each arm's own correct queries: theorem's are generated queries that scored 1.0,
Cypher's are CypherBench's gold. Outcomes are \textbf{rejected} (an error instead
of an answer), \textbf{empty} (ran, returned nothing, indistinguishable from a
true empty answer), \textbf{wrong} (ran, returned rows that are not the right
rows), and \textbf{inert} (the mutation did not change the answer; not counted).

\begin{center}
\begin{tabular}{lrrrrr}
\toprule
arm & mutants & rejected & empty & wrong & undetectable \\
\midrule
theorem     & 1{,}928 & \textbf{1{,}811} & 69 & 48 & \textbf{6.1\%} \\
text2cypher & 1{,}997 & \textbf{0} & 1{,}281 & 716 & \textbf{100\%} \\
\bottomrule
\end{tabular}
\end{center}

Every broken Cypher query returned rows or a confident empty set. Not one
produced an error. Three caveats ship on the page rather than buried:

\begin{itemize}
  \item \textbf{theorem's 6.1\% is entirely direction}, and only on edges whose
    two roles hold the same class. \texttt{hasFather(subj: person, obj: person)}
    is type-correct either way round, so swapping the roles asks a different
    question rather than an invalid one, and no schema check can tell. On
    \texttt{nba}, which has no such edge, theorem catches everything: 0.0\%
    across 843 mutants. On \texttt{fictional\_character}, which has five, half
    the direction mutants survive.
  \item \textbf{Neo4j does notify the driver.} A missing label, relationship type
    or property raises a \texttt{01N42} notification alongside a successful
    result: 1{,}532 of the 1{,}997. It is a warning on a call that succeeded, it
    never fires on a reversed arrow, and no published text2cypher pipeline reads
    it. It is counted in its own column.
  \item \textbf{A mutation is not a model.} This measures what a language does
    with a broken query, not how often a model breaks one.
\end{itemize}

\subsection{Frontier model (new)}

The standing objection to a new query language is that it is scaffolding for
weak models, and that a better model writes Cypher well enough to make the
exercise pointless. It predicts the gap narrows with model scale.

498 questions sampled from the test set, seeded and stratified by graph and
match category, run through both languages on both models, so the comparison is
paired across models as well as across languages.

\begin{center}
\begin{tabular}{lrrr}
\toprule
Model & theorem & text2cypher & gap \\
\midrule
Haiku 4.5 & 84.3\% & 75.5\% & $\mathbf{+8.8}$ \\
Sonnet 5  & 74.1\% & 64.3\% & $\mathbf{+9.8}$ \\
\bottomrule
\end{tabular}
\end{center}

On Sonnet 5, theorem alone answers 96 of these questions and text2cypher alone
answers 47. Exact McNemar $p = 0.0001$. \textbf{The gap does not close; it is
slightly wider on the frontier model.}

The same run found something nobody was looking for: \emph{both} arms score about
ten points worse on the frontier model than on the small one, on identical
questions, by roughly the same amount in each language. That makes it a property
of the task rather than of either language. The obvious guess is that a benchmark
rewarding terse literal translation penalises a model that elaborates, but
nothing in this run measures the mechanism, and it is published as a question
rather than an answer.

\subsection{Prompt cost (new)}

theorem carries a tutorial in every prompt because the model has never seen the
language; Cypher carries none because it has. That is a fixed cost. Against it,
theorem's schema render is much cheaper per class than the JSON schema a
text2cypher prompt sends, so the two costs are lines with different slopes.

The roadmap had quoted a crossover at 18--20 classes, fitted over graphs with
9--13 classes, which is both theorem's worst region and the region least able to
show a crossing. Measured instead, on the seven test schemas and on their union
(a real 40-class, 62-edge-type schema assembled from real ones):

\begin{itemize}
  \item theorem: \textbf{39 tokens per class}, on 1{,}357 tokens of fixed tutorial.
  \item text2cypher: \textbf{85 tokens per class}, on 99.
  \item The lines cross at \textbf{31 classes}.
\end{itemize}

The crossing is inside the measured range rather than past it: on the union,
theorem's prompt is the smaller of the two, 2{,}793 tokens against 3{,}186.

\section{The guards}

Published results are a property of a version of the code, and an engine change
can move them without anyone noticing. Two guards exist and both earned their
keep.

\paragraph{Replay.} \texttt{eval/verify\_replay.py} re-executes the exact queries
that were scored, against the exact stores they were scored on, and reports any
question whose score moved. It generates nothing and calls no model, so it costs
minutes rather than hours. Every engine change this session was checked this way:
\textbf{2{,}348 questions, zero moved} for every change except one deliberate
language widening, which moved nine questions, \emph{all of them from 0 to 1}.

It also resolved a provenance problem. The final execution phase resumed from
partial files scored before the engine changes landed, which would have made the
published number a blend of two engines. Replaying all 2{,}348 on the finished
engine reproduced every score, so the blend was not a blend.

\paragraph{Prompt fingerprint.} Frozen query files are named by the hash of the
tutorial that produced them, and each scored graph now writes its own fingerprint
beside its results. A report reads the hash from the runs rather than recomputing
it, because recomputing prints whatever tutorial happens to be checked out. That
is precisely how a report comes to claim a prompt it never ran, and the
CypherBench page was doing it.

\section{What the sampling found}
\label{sec:sampling}

Partway through regeneration, sampling the queries the new prompt was producing
showed it verifying \textbf{93.8\%} where the old prompt verified \textbf{97.8\%}
on identical questions. A four-point regression, invisible in the agent loop
where a repair turn recovers it, and fatal to a one-shot benchmark.

Almost the entire gap was one rejected spelling:

\begin{lstlisting}[style=err]
follow c locatedIn lake as l where l.area_km2 < 390000   # rejected
follow c locatedIn lake where area_km2 < 390000 as l     # accepted
\end{lstlisting}

They mean the same thing. A follow's condition is about the node being arrived
at, and \texttt{l} is the name for that node. Models write the qualified form
constantly because it is the unambiguous one.

The parser now drops a leading binding name from a condition when it is the name
that statement itself creates. Verification went to \textbf{97.9\%}, level with
the long prompt, \emph{without touching the prompt}, which would have invalidated
a thousand cached generations. That fix is why the headline held at $-0.04$
points instead of dropping four.

Two residual failure classes remain, seven questions in about 330, both recorded
in the roadmap with the evidence rather than quietly fixed: comparing two
bindings' properties (\texttt{keep c where birth\_loc.name = death\_loc.name}),
which the language cannot express, and role guessing on edges whose roles are
named for their classes, which the error already teaches and a repair turn
already fixes.

\section{Defects found in shipped code}

Eight, all predating the session, all fixed and covered by tests.

\begin{center}
\small
\begin{tabular}{p{0.34\textwidth}p{0.58\textwidth}}
\toprule
Defect & Symptom \\
\midrule
\texttt{upto any} walked paths, not nodes & A 60-node ring did not finish in 400\,s for a 60-row answer \\
Two writers on one store & Both allocated \texttt{\#e-1}, both wrote, one node survived, \textbf{no error} \\
\texttt{snapshot()} had no callers & The log grew forever; every startup replayed all history \\
Reads were writes & 4{,}178 log records and 129\,ms for one ordinary query; 17\,ms and none after \\
\texttt{or} and \texttt{keep} unreachable & Missing from the session's read-statement list, so any program using them was routed to the write path and rejected \\
Name reuse ignored under \texttt{upto} & The join rule held or not depending on whether you wrote \texttt{upto} \\
\texttt{find} ignored subclasses & \texttt{derive class} partitioned instances away from every query against the base; the shipped ingest pipeline answered nothing for \texttt{find piece} on a store full of pieces \\
``Partial execution never happens'' & True of verification, false of execution: four asserts against a quota of two wrote two nodes and printed only an error \\
\bottomrule
\end{tabular}
\end{center}

The reads-as-writes defect deserves a note. All four benchmark harnesses were
monkey-patching \texttt{store.apply} on the read path, with a comment explaining
why, which meant the published latency described an engine that was not shipping.
The patch is gone from all four; they now measure what ships.

\section{Defects found in the benchmark harness}

Four, all caught before publication.

\begin{enumerate}
  \item \textbf{Generation aborted on one transport failure.} A sustained rate
    limit failed 524 of 2{,}348 calls, and a single unrecovered failure would
    have taken the other 2{,}347 with it hours in. Generation now records the
    miss and continues, and \emph{refuses to write a frozen query file with holes
    in it}, because holes score as wrong answers. When the throttle hit, that is
    exactly what happened: no frozen file, and the execution phase reported
    \texttt{INCOMPLETE} rather than publishing a partial number.
  \item \textbf{Retry patience did not match the failure.} Five attempts spanning
    75 seconds rides out a transport blip, not a rate limit. Now eight attempts
    reaching about eleven minutes, recording stdout as well as stderr, because
    the CLI reports usage limits on stdout and \texttt{rc=1} with nothing after
    it is not a diagnosis.
  \item \textbf{The frontier Cypher arm got the wrong prompt.} It was handed the
    raw schema JSON instead of the benchmark's own schema rendering. The model
    dutifully wrote \texttt{\{label: "..."\}} for every node, matched nothing,
    and the arm \textbf{scored 0.000}.
  \item \textbf{Then it got the wrong scoring.} Neo4j records were reshaped by a
    helper meant for gold rows. The arm \textbf{scored 0.000 again.}
\end{enumerate}

Both frontier bugs would have published ``theorem 74\% against Cypher 0\%''. They
were absurd enough to notice. A subtler version would not have been, which is the
argument for the replay checker existing at all. A fifth, smaller one: the report
generator read the new provenance files as result files, because they matched the
same glob.

\section{Production readiness}

Benchmarks are half of what ``publishable'' needs; the other half is an engine
somebody can run. Beyond the eight defects above:

\begin{itemize}
  \item \textbf{One writer per store}, enforced by an advisory lock that a crash
    releases, refusing the second opener by name.
  \item \textbf{Automatic WAL compaction}, amortized so the threshold is at least
    the size of the live data. Without that, a million-node ingest rewrites a
    million-record run file on every threshold crossing.
  \item \textbf{A row and wall-clock ceiling on every read}, in the engine rather
    than in each caller.
  \item \textbf{A stated durability contract}: a committed record survives the
    process dying, because the write has reached the operating system. It does
    not survive power loss, because the write path does not fsync. Snapshots do.
  \item \textbf{\texttt{theorem load}} for CSV and JSONL, refusing the file rather
    than writing half of it.
  \item \textbf{An import surface.} \texttt{import theorem} previously returned a
    \texttt{hello()} stub from the project template.
  \item \textbf{The prompt and the agent loop ship.} They lived in \texttt{eval/},
    so no user had the prompt the numbers describe. They are
    \texttt{theorem.prompt} and \texttt{theorem.answer} now, byte-identical, and
    the harness re-exports them.
  \item \textbf{A property index}, built on first use above 5{,}000 nodes:
    299\,ms to 20\,ms for a common value, 341\,ms to 0.8\,ms for a rare one, and
    304\,ms to nothing for a value no node has.
  \item \textbf{\texttt{theorem.canonical}}, \texttt{theorem stats},
    \texttt{Store.refresh()} so a reader can follow a writer, and a release
    workflow that refuses to publish a tag disagreeing with the package version.
\end{itemize}

540 tests on Python 3.11, 3.12 and 3.13. Lint clean and pinned. Coverage 89.8\%.
The wheel builds and imports in an isolated environment.

\section{What the numbers do not support}

Stated plainly, because someone will look for these.

\begin{itemize}
  \item \textbf{The agent loop is a tie.} 87.9\% against 85.0\% at $n=240$ is
    $p = 0.392$. The first-attempt gap is real and large; the third-attempt gap
    is not established.
  \item \textbf{The two arms have different prompts.} theorem's is a tutorial
    with an EBNF grammar and nine worked examples; Cypher's is the official
    zero-shot prompt. This is each system with its natural prompting, not a
    matched-prompt comparison, so the 7.6 points mixes the language with how it
    is taught. \emph{This is the strongest available criticism, and answering it
    means running the Cypher arm few-shot.}
  \item \textbf{\texttt{nba} is in theorem's tuning set.} Hence the 76.85\%
    held-out figure, which is the number to quote under scrutiny.
  \item \textbf{Nothing has been run against a 40-class graph.} The prompt-cost
    crossover is measured; accuracy at that scale is not.
  \item \textbf{The frontier drop is unexplained.} Both languages lose about ten
    points on the larger model, and the mechanism is a guess.
  \item \textbf{Writes are not transactional.} A program that verifies can still
    fail while running, and each write commits as it happens.
\end{itemize}

\section{Reproducing any of it}

Per-question results for both arms, both frozen query files, and each graph's
prompt fingerprint are committed under \texttt{eval/out/public/}. Every table
above can be recomputed rather than trusted.

\begin{lstlisting}[style=gl]
uv run python -m eval.run_public all      # CypherBench, both phases
uv run python -m eval.run_cypher_public all   # the Cypher control (needs docker)
uv run python -m eval.run_agent run --graph soccer --n 120 --arm both
uv run python -m eval.run_silent --graph nba  # silent failure
uv run python -m eval.token_crossover --report # prompt cost, no model calls
uv run python -m eval.run_frontier gen --model claude-sonnet-5
uv run python -m eval.verify_replay       # did an engine change move a score?
\end{lstlisting}

\section{The short version}

\begin{center}
\begin{tabular}{lrr}
\toprule
 & theorem & text2cypher \\
\midrule
CypherBench, 2{,}348 questions      & \textbf{77.98\%} & 70.36\% \\
Agent loop, $n=240$, solve@3        & 87.9\% & 85.0\% \\
Agent loop, $n=240$, solve@1        & \textbf{82.9\%} & 72.1\% \\
Silent failure, 3{,}925 mutants     & refuses \textbf{1{,}811/1{,}928} & refuses \textbf{0/1{,}997} \\
Frontier model, 498 questions       & \textbf{74.1\%} & 64.3\% \\
\bottomrule
\end{tabular}
\end{center}

\vspace{0.6em}

The headline survived a prompt change that should have cost it four points,
because sampling caught the regression and a language widening closed it. The
strongest objection to the project, that model scale dissolves the advantage, was
tested directly and does not hold. The claim the project is actually built on,
that a wrong query is loud rather than plausible, now has a number: 6.1\% against
100\%.

And the number the repository published at the start of the day could not be
reproduced by the code that shipped it, which is the finding that made the rest
of the day necessary.

\end{document}
